Un haut-parleur, qu'on peut considérer comme une source ponctuelle qui émet uniformément dans toutes les directions, a une puissance acoustique de \qty{\v{P}}{W}. On néglige les réflexions et l'absorption par l'air.

\begin{enumerate}[label=\alph*)]
    \item Déterminez l'intensité sonore à \qty{\v{r}}{m} du haut-parleur.
    \item Déterminez le niveau sonore à cette distance.
    \item À quelle distance du haut-parleur le niveau sonore n'est-il plus que de \qty{\v{b2}}{dB} ?
\end{enumerate}
